% Claim statement file for row claim_3_2 -- "AcyclicIffTopologicalOrder".
% Auto-generated stub by scaffold/solve_chapter.py at row start. The
% manager agent (or a worker it dispatches) fills the body below with the
% claim's statement(s) from the lecture notes -- statement only, no proof
% (proofs live in the sibling `claim_3_2_proof_*.tex` file).
%
% This file is a sibling of `main.tex` inside `<subsection>/tex/`. The
% `subfiles` package lets it render both standalone and as part of
% `main.tex`. Note: `main.tex` does NOT \subfile this statement file -- the
% sibling `_proof_` file restates the statement above the proof, so
% including both in `main.tex` would be redundant. This file exists for
% standalone reading / grep convenience.

\documentclass[main]{subfiles}

\begin{document}
% ref: claim_3_2 (statement)
% title: AcyclicIffTopologicalOrder
\def\rowref{claim\_3\_2}
\phantomsection\label{claim_3_2}

\begin{Lem}[AcyclicIffTopologicalOrder]
    Let $G = (J, V, E, L)$ be a CDMG (def \ref{def-cdmg}). Then $G$ is acyclic (in the sense of def \ref{def-acylic}) if and only if there exists a topological order of $G$ (in the sense of def \ref{def-topological-order}); equivalently, $G$ is acyclic if and only if there exists a strict total order $<$ on $J \cup V$ -- i.e.\ a binary relation $<$ on $J \cup V$ that is irreflexive, transitive, and trichotomous -- such that, for every $v, w \in J \cup V$, $v \in \Pa^G(w) \implies v < w$ (using the parent-set notation $\Pa^G$ of def \ref{def_3_5}, item~i.). The label set $L$ plays no role on either side of the biconditional.
\end{Lem}

\end{document}
